% Preamble template for Tufte-style applied econometrics lecture notes.
% Do not compile this file directly; compile the main notes file instead.
\documentclass{tufte-handout}

\usepackage[T1]{fontenc}
\usepackage[utf8]{inputenc}
\usepackage{ntheorem}
\usepackage{graphicx}
\usepackage{amsmath}
\usepackage{amssymb}
\usepackage{mathtools}
\usepackage{hyperref}
\usepackage{epigraph}
\usepackage{booktabs}
\usepackage{array}
\theoremstyle{break}

\hypersetup{
  colorlinks,
  urlcolor = blue,
  pdfauthor={Swapnil Singh},
  pdfkeywords={econometrics, causal inference, applied econometrics},
  pdftitle={Lecture Notes for Applied Econometrics},
  pdfpagemode=UseNone
}

\newtheorem{ruleN}{Rule}
\newtheorem{thmN}{Theorem}
\newtheorem{assN}{Assumption}
\newtheorem{defN}{Definition}
\newtheorem{exmp}{Example}
\newtheorem{cmt}{Comment}
\newtheorem{proof}{Proof}

\newcommand{\continuation}{??}
\newtheorem*{excont}{Example \continuation}
\newenvironment{continueexample}[1]
 {\renewcommand{\continuation}{\ref{#1}}\excont[continued]}
 {\endexcont}

\newcommand{\bY}{\mathbf{Y}}
\newcommand{\bX}{\mathbf{X}}
\newcommand{\bD}{\mathbf{D}}
\newcommand{\E}{\mathbb{E}}
\newcommand{\Prob}{\mathbb{P}}
\newcommand{\Var}{\mathrm{Var}}
\newcommand{\Cov}{\mathrm{Cov}}
\newcommand{\se}{\mathrm{SE}}
\newcommand{\CI}{\mathrm{CI}}
\newcommand{\ATE}{\mathrm{ATE}}
\newcommand{\ATT}{\mathrm{ATT}}
\newcommand{\TOT}{\mathrm{TOT}}
\newcommand{\LATE}{\mathrm{LATE}}
\newcommand{\ITT}{\mathrm{ITT}}
\newcommand{\RD}{\mathrm{RD}}
\newcommand{\DiD}{\mathrm{DiD}}
\newcommand{\MMSE}{\mathrm{MMSE}}
\newcommand{\OLS}{\mathrm{OLS}}
\newcommand{\Normal}{\mathcal{N}}
\newcommand{\plim}{\operatorname*{plim}}
\DeclareMathOperator*{\argmin}{arg\,min}
\DeclareMathOperator{\Supp}{Supp}

\newcommand\independent{\protect\mathpalette{\protect\independenT}{\perp}}
\def\independenT#1#2{\mathrel{\rlap{$#1#2$}\mkern2mu{#1#2}}}
\newcommand{\indep}{\perp\!\!\!\perp}

\usepackage{cleveref}
\crefname{appsec}{appendix}{appendices}
\crefname{appsubsec}{appendix}{appendices}
\crefname{assumption}{assumption}{assumptions}
\crefname{equation}{equation}{equations}
\crefname{exmp}{example}{examples}
\crefname{assN}{assumption}{assumptions}
\crefname{cmt}{comment}{comments}
\crefname{defN}{definition}{definitions}

\usepackage[nolist]{acronym}
\begin{acronym}
  \acro{ATE}{average treatment effect}
  \acro{ATT}{average treatment effect on the treated}
  \acro{CIA}{conditional independence assumption}
  \acro{CI}{confidence interval}
  \acro{CLT}{central limit theorem}
  \acro{DiD}{difference-in-differences}
  \acro{FE}{fixed effects}
  \acro{FWL}{Frisch-Waugh-Lovell}
  \acro{IV}{instrumental variables}
  \acro{LATE}{local average treatment effect}
  \acro{OLS}{ordinary least squares}
  \acro{OVB}{omitted-variable bias}
  \acro{PO}{potential outcomes}
  \acro{RCT}{randomized controlled trial}
  \acro{RD}{regression discontinuity}
  \acro{SUTVA}{stable unit treatment value assignment}
  \acro{TWFE}{two-way fixed effects}
\end{acronym}

\def\inprobHIGH{\,{\buildrel p \over \rightarrow}\,}
\def\inprob{\,{\inprobHIGH}\,}
\def\indistHIGH{\,{\buildrel d \over \rightarrow}\,}
\def\indist{\,{\indistHIGH}\,}

\usepackage[many]{tcolorbox}
\definecolor{main}{HTML}{3F6EA7}
\definecolor{sub}{HTML}{F2F7FD}
\definecolor{sub2}{HTML}{FFF5E8}
\definecolor{mainline}{HTML}{D3DFEE}
\definecolor{warn}{HTML}{C7771B}
\definecolor{warnline}{HTML}{EAD9C0}

\tcbset{
    enhanced,
    breakable,
    colback = white,
    boxrule = 0.5pt,
    arc = 2pt,
    left = 9pt,
    right = 9pt,
    top = 7pt,
    bottom = 7pt,
    before skip = 0.3cm,
    after skip = 0.45cm
}

\newtcolorbox{boxD}{
    colback = sub,
    colframe = mainline,
    borderline west = {2.2pt}{0pt}{main},
    borderline north = {0.6pt}{0pt}{main!25},
    borderline south = {0.6pt}{0pt}{main!15}
}

\newtcolorbox{boxF}{
    colback = sub2,
    colframe = warnline,
    borderline west = {2.2pt}{0pt}{warn},
    borderline north = {0.6pt}{0pt}{warn!25},
    borderline south = {0.6pt}{0pt}{warn!15}
}

\usepackage{colortbl}
\usepackage{tikz}
\usepackage{verbatim}
\usetikzlibrary{positioning}
\usetikzlibrary{snakes}
\usetikzlibrary{calc}
\usetikzlibrary{arrows}
\usetikzlibrary{decorations.markings}
\usetikzlibrary{shapes.misc}
\usetikzlibrary{matrix,shapes,arrows,fit,tikzmark}


\title{Applied Econometrics: Session 2 Notes}
\author{Swapnil Singh}
\date{}

\hypersetup{
  pdftitle={Applied Econometrics: Session 2 Notes},
  pdfauthor={Swapnil Singh},
  pdfkeywords={regression, CEF, selection on observables, matching, propensity score, bad controls}
}

\setlength{\headheight}{26pt}

\begin{document}
\maketitle

\begin{abstract}
This lecture explains what regression means before we use it for causality.
Regression is a linear approximation to conditional means. It becomes causal
only after an identification argument connects observed conditional means to
potential-outcome means. We study the conditional expectation function, the law
of iterated expectations, regression anatomy, conditional independence,
omitted-variable bias, bad controls, matching, propensity scores, inverse
probability weighting, and practical issues such as grouped data, weights, and
limited dependent variables.
\end{abstract}

\begin{boxD}
\textbf{Reading and objective.} These notes accompany Session 2, which builds
on the Mostly Harmless Econometrics treatment of regression, controls, matching,
and selection on observables. The objective is to understand what regression
always estimates mechanically, and then to state the additional assumptions
needed before a regression coefficient can be read as a causal effect.
\end{boxD}

\section{Regression Starts with Conditional Means}

Lecture 1 introduced the causal problem: we want missing potential outcomes.
Lecture 2 asks how regression can help. The answer has two layers.
\begin{enumerate}
    \item Mechanically, regression summarizes conditional means.
    \item Causally, regression is useful only if conditioning creates credible
    comparisons.
\end{enumerate}

\begin{boxF}
\textbf{Core warning.} Regression is not a causal machine. It is a statistical
summary. Causal interpretation comes from assumptions about treatment assignment,
controls, timing, and comparison groups.
\end{boxF}

\subsection{The Conditional Expectation Function}

For outcome \(Y_i\) and covariates \(X_i\), the conditional expectation function
is
\[
    m(x)=\E[Y_i\mid X_i=x].
\]
It gives the average outcome among units with covariate value \(x\). If \(X_i\)
is schooling, \(m(12)\) is the average outcome among people with 12 years of
schooling. If \(X_i\) is a vector, \(m(x)\) is an average within a covariate
cell or covariate profile.

\begin{defN}[Conditional expectation function]
The CEF of \(Y_i\) given \(X_i\) is the function
\[
    m(X_i)=\E[Y_i\mid X_i].
\]
It is the best unrestricted mean-square predictor of \(Y_i\) using \(X_i\).
\end{defN}

The CEF is descriptive unless we add causal assumptions. For example,
\(\E[\log w_i\mid S_i=16]\) is the average log wage of workers with 16 years of
schooling. It does not by itself say what the same workers would have earned if
they had received less schooling.

\subsection{Law of Iterated Expectations}

The law of iterated expectations says
\[
    \E[Y_i]=\E\{\E[Y_i\mid X_i]\}.
\]
First average \(Y_i\) within covariate groups. Then average those conditional
means over the distribution of \(X_i\). This identity is the bridge from
cell-specific objects to population averages.

\begin{exmp}[Two schooling groups]
Suppose 40 percent of workers have 12 years of schooling and average wage 20,
while 60 percent have 16 years of schooling and average wage 35. The overall
average wage is
\[
    0.4(20)+0.6(35)=29.
\]
This is exactly the law of iterated expectations in a discrete example.
\end{exmp}

\subsection{CEF Decomposition}

Write
\[
    Y_i = \E[Y_i\mid X_i] + \varepsilon_i.
\]
Then
\[
    \E[\varepsilon_i\mid X_i]=0.
\]
This says that once we subtract the conditional mean, the remaining error has
zero average within every covariate group. Therefore the residual is uncorrelated
with any function of \(X_i\).

\begin{thmN}[CEF decomposition]
Let \(\varepsilon_i=Y_i-\E[Y_i\mid X_i]\). Then
\[
    \E[\varepsilon_i\mid X_i]=0
\]
and for any function \(g\),
\[
    \E[g(X_i)\varepsilon_i]=0.
\]
\end{thmN}

\begin{proof}
The first statement follows from
\[
    \E[\varepsilon_i\mid X_i]
    =
    \E[Y_i-\E[Y_i\mid X_i]\mid X_i]
    =
    \E[Y_i\mid X_i]-\E[Y_i\mid X_i]
    =
    0.
\]
For the second statement, use iterated expectations:
\[
    \E[g(X_i)\varepsilon_i]
    =
    \E\{g(X_i)\E[\varepsilon_i\mid X_i]\}=0.
\]
\end{proof}

\section{Population Regression}

Regression is a linear approximation problem. Let \(X_i\) include an intercept.
The population linear regression coefficient is
\[
    \beta
    =
    \argmin_b \E[(Y_i-X_i'b)^2].
\]
The first-order condition is
\[
    \E[X_i(Y_i-X_i'\beta)]=0.
\]
If \(\E[X_iX_i']\) is invertible, then
\[
    \beta = \E[X_iX_i']^{-1}\E[X_iY_i].
\]

\begin{defN}[Population regression coefficient]
The population regression coefficient is the coefficient vector of the best
linear predictor of \(Y_i\) from \(X_i\) in mean squared error.
\end{defN}

The residual \(e_i=Y_i-X_i'\beta\) is orthogonal to the regressors:
\[
    \E[X_i e_i]=0.
\]
This orthogonality is an algebraic property of least squares. It does not mean
the regressors are exogenous in the causal sense.

\begin{boxF}
\textbf{Important distinction.} Regression residuals are orthogonal to included
regressors by construction. That is not the same as saying the regressors are
as-good-as-random.
\end{boxF}

\subsection{Regression Anatomy}

In a simple regression of \(Y_i\) on an intercept and scalar \(x_i\),
\[
    \beta_1=\frac{\Cov(Y_i,x_i)}{\Var(x_i)}.
\]
In multiple regression, the coefficient on \(x_{ki}\) has the same form after
removing the part of \(x_{ki}\) explained by the other regressors. Let
\(\tilde{x}_{ki}\) be the residual from regressing \(x_{ki}\) on the other
covariates. Then
\[
    \beta_k
    =
    \frac{\Cov(Y_i,\tilde{x}_{ki})}{\Var(\tilde{x}_{ki})}.
\]

\begin{thmN}[Regression anatomy]
A multiple-regression coefficient is a bivariate slope using only the residual
variation in the regressor that remains after partialling out the other
covariates.
\end{thmN}

\begin{boxD}
\textbf{Plain-language interpretation.} The coefficient on schooling in a wage
regression with age and region controls compares workers who differ in schooling
after removing the part of schooling predicted by age and region. This is why
control choice matters so much.
\end{boxD}

\subsection{Saturated Models}

If all controls are discrete and we include a dummy for every cell, regression
is saturated. A saturated regression reproduces cell means exactly. If we leave
out interactions, we impose restrictions. For example, controlling for gender
and region without a gender-by-region interaction assumes the gender difference
does not vary by region, unless other variables capture it.

\begin{boxF}
\textbf{Student warning.} ``Controlling for \(X\)'' is vague. A linear control,
a set of dummies, and a fully interacted saturated model are different
statistical objects.
\end{boxF}

\section{Regression and Causality}

Regression becomes causal when the conditional comparisons in the data reveal
potential-outcome comparisons. The key assumption in this lecture is selection
on observables, also called conditional independence.

\begin{assN}[Conditional independence]
For binary treatment \(D_i\) and covariates \(X_i\),
\[
    (Y_{1i},Y_{0i}) \indep D_i \mid X_i.
\]
Conditional on \(X_i\), treatment is as good as randomly assigned.
\end{assN}

If conditional independence holds, then for each \(x\),
\[
    \E[Y_i\mid D_i=1,X_i=x]=\E[Y_{1i}\mid X_i=x],
\]
and
\[
    \E[Y_i\mid D_i=0,X_i=x]=\E[Y_{0i}\mid X_i=x].
\]
Therefore the conditional observed difference
\[
    \E[Y_i\mid D_i=1,X_i=x]-\E[Y_i\mid D_i=0,X_i=x]
\]
identifies the treatment effect at covariate value \(x\):
\[
    \E[Y_{1i}-Y_{0i}\mid X_i=x].
\]

\begin{boxD}
\textbf{What CIA buys.} It lets us compare treated and untreated units within
the same covariate cell as if treatment were randomized inside that cell.
\end{boxD}

\subsection{From Cell Effects to Population Effects}

After estimating \(X\)-specific effects, we must decide how to average them.
The ATE averages using the population distribution of \(X\):
\[
    \ATE = \E\left[\E[Y_{1i}-Y_{0i}\mid X_i]\right].
\]
The ATT averages using the distribution of \(X\) among treated units:
\[
    \ATT = \E\left[\E[Y_{1i}-Y_{0i}\mid X_i]\mid D_i=1\right].
\]
These differ when treatment effects vary and treated units have a different
covariate distribution than the full population.

\begin{exmp}[Training program]
Suppose a training program helps low-experience workers by 8 units and
high-experience workers by 2 units. If most treated workers are low-experience,
the ATT can be much larger than the ATE. Reporting ``the'' treatment effect
without saying how effects are averaged is incomplete.
\end{exmp}

\section{Omitted Variables and Bad Controls}

Consider the causal model
\[
    Y_i=\alpha+\tau D_i+\gamma Z_i+u_i,
\]
but suppose we estimate
\[
    Y_i=a+bD_i+v_i.
\]
Then the short-regression coefficient is
\[
    b=\tau+\gamma\frac{\Cov(D_i,Z_i)}{\Var(D_i)}.
\]

\begin{thmN}[Omitted-variable bias]
When the true linear model includes \(Z_i\) but \(Z_i\) is omitted, the
coefficient on \(D_i\) equals the causal effect plus
\[
    \gamma\delta,
    \qquad
    \delta=\frac{\Cov(D_i,Z_i)}{\Var(D_i)}.
\]
Bias requires two ingredients: the omitted variable predicts the outcome and is
correlated with treatment.
\end{thmN}

\begin{boxD}
\textbf{Sign rule.} Bias is positive if the omitted variable raises the outcome
and treated units have more of it. Bias is negative if the omitted variable
raises the outcome and treated units have less of it. Always write both signs.
\end{boxD}

\begin{exmp}[Schooling and wages]
If ability raises wages and high-ability workers obtain more schooling, omitting
ability biases the schooling coefficient upward. If a scholarship program
targets disadvantaged students who would otherwise have low earnings, omitting
disadvantage may bias the program coefficient downward.
\end{exmp}

\subsection{Good Controls}

A good control is usually a pre-treatment variable that affects both treatment
assignment and potential outcomes. It helps compare treated and untreated units
that were similar before treatment.

\begin{defN}[Good control]
A good control is determined before treatment and blocks a back-door path
between treatment and outcome without itself being caused by treatment.
\end{defN}

\subsection{Bad Controls}

Bad controls are variables that make the comparison worse. The most common bad
control is a mediator, a variable caused by treatment on the path to the
outcome. Controlling for it removes part of the effect we may want to estimate.

\begin{defN}[Mediator]
A mediator is a variable \(M_i\) affected by treatment that in turn affects the
outcome. If the target is the total effect of treatment, controlling for
\(M_i\) blocks part of the causal channel.
\end{defN}

Another bad-control problem is conditioning on a collider or selected sample.
If a variable is affected by treatment and by unobserved determinants of the
outcome, conditioning on it can create artificial association.

\begin{boxF}
\textbf{Control-selection rule.} Controls are not chosen because they make
\(R^2\) larger. They are chosen because the causal design says they are needed
and safe.
\end{boxF}

\section{Matching}

Matching implements the conditional-independence idea directly. For each treated
unit, find untreated units with similar covariates and use their outcomes as
counterfactuals.

\begin{defN}[Matching estimator, ATT intuition]
For treated unit \(i\), let \(j(i)\) be a matched untreated unit with similar
covariates. A simple matching estimator for the ATT is
\[
    \widehat{\ATT}
    =
    \frac{1}{N_1}\sum_{i:D_i=1}\left(Y_i-Y_{j(i)}\right).
\]
\end{defN}

Exact matching is transparent but quickly becomes impossible when \(X_i\) has
many components. Nearest-neighbor matching is feasible but introduces judgment:
how close is close enough? Which distance metric should be used? Should controls
be reused? Should bad matches be dropped?

\begin{boxF}
\textbf{Curse of dimensionality.} With many covariates, exact matches become
rare. Matching then depends heavily on distance metrics and overlap.
\end{boxF}

\subsection{Common Support}

Common support means treated and untreated units overlap in covariate space.
If treated units have covariate values that never occur among controls, the data
contain no credible comparison for those treated units.

\begin{assN}[Overlap]
For treatment probability \(p(X_i)=\Pr(D_i=1\mid X_i)\),
\[
    0 < p(X_i) < 1
\]
for covariate values in the target population.
\end{assN}

\section{Propensity Scores}

The propensity score is
\[
    p(X_i)=\Pr(D_i=1\mid X_i).
\]
It compresses many covariates into one scalar treatment probability.

\begin{thmN}[Propensity-score theorem]
If
\[
    (Y_{1i},Y_{0i})\indep D_i\mid X_i,
\]
then
\[
    (Y_{1i},Y_{0i})\indep D_i\mid p(X_i).
\]
\end{thmN}

The theorem says that if conditioning on all covariates is enough, then
conditioning on the propensity score is also enough. It does not say that a
poorly estimated propensity score solves hidden confounding. It also does not
remove the need for overlap.

\begin{boxF}
\textbf{Beginner trap.} The goal of a propensity-score model is not maximum
classification accuracy. A model that predicts treatment almost perfectly may
destroy overlap. The goal is balance between treated and untreated units.
\end{boxF}

\subsection{Inverse Probability Weighting}

Inverse probability weighting reweights observations to construct a pseudo-
population in which treatment is independent of covariates. For the ATE,
\[
    \widehat{\ATE}_{IPW}
    =
    \frac{1}{n}\sum_{i=1}^n
    \left[
    \frac{D_iY_i}{\hat{p}(X_i)}
    -
    \frac{(1-D_i)Y_i}{1-\hat{p}(X_i)}
    \right].
\]

The intuition is simple. Treated units who were unlikely to be treated get large
weight because they represent many similar untreated-looking units. Control
units who were unlikely to be controls also get large weight. Extreme weights
are therefore a warning sign about overlap.

\begin{boxD}
\textbf{Family resemblance.} Regression, matching, propensity-score matching,
and IPW all try to compare treated and untreated units with similar observed
covariates. They differ in modeling choices and weighting, not in the basic
identification assumption.
\end{boxD}

\section{Practical Regression Issues}

\subsection{Weights and Grouped Data}

If data are grouped into cells, a regression on cell means should usually weight
by cell size when the target is the microdata regression. A cell with 1,000
observations contains more information about the micro population than a cell
with 5 observations.

\begin{thmN}[Grouped regression weighting]
When cell means are used in place of microdata and the target is the original
microdata least-squares coefficient, weighted least squares with weights equal
to cell counts recovers the microdata regression under the corresponding
cell-mean specification.
\end{thmN}

\begin{boxF}
\textbf{Heteroskedasticity warning.} Do not weight mechanically just because
errors are heteroskedastic. Robust standard errors often handle inference
without changing the estimand. Weighting changes the coefficient by changing
which observations matter more.
\end{boxF}

\subsection{Limited Dependent Variables}

OLS can be useful even when the outcome is binary. If \(Y_i\in\{0,1\}\), then
\[
    \E[Y_i\mid X_i]=\Pr(Y_i=1\mid X_i).
\]
A linear probability model approximates this conditional probability. It may
predict values outside \([0,1]\), but the coefficient can still be a useful
average marginal effect.

Probit and logit impose nonlinear probability curves. They are often used for
prediction or because the nonlinear shape is substantively helpful. But their
coefficients are not directly comparable to OLS coefficients; marginal effects
are the relevant comparison object.

\subsection{Conditional-on-Positive Outcomes}

Sometimes researchers study outcomes only when they are positive, such as wages
only among employed workers or medical spending only among users. This can be a
bad-control problem because treatment may affect whether the outcome is
positive.

\begin{boxF}
\textbf{Conditional-on-positive warning.} If treatment changes the probability
that \(Y_i>0\), then analyzing only observations with \(Y_i>0\) conditions on a
post-treatment event. The result is generally not the total causal effect.
\end{boxF}

\section{Worked Problems}

\begin{exmp}[Using the law of iterated expectations]
Suppose half the sample is young with average outcome 10, and half is old with
average outcome 16. Then
\[
    \E[Y]=0.5(10)+0.5(16)=13.
\]
The unconditional mean is a weighted average of conditional means.
\end{exmp}

\begin{exmp}[Omitted-variable bias sign]
Let \(D_i\) be college attendance and \(Z_i\) be ability. Suppose ability raises
wages, so \(\gamma>0\), and ability is positively correlated with college, so
\(\delta>0\). The omitted-variable bias is \(\gamma\delta>0\). A regression of
wages on college alone overstates the causal return to college.
\end{exmp}

\begin{exmp}[ATT versus ATE weights]
There are two groups. Low-experience workers have treatment effect 8 and make up
70 percent of treated workers but 40 percent of the population. High-experience
workers have treatment effect 2. Then
\[
    \ATT=0.7(8)+0.3(2)=6.2,
\]
while
\[
    \ATE=0.4(8)+0.6(2)=4.4.
\]
The estimand changes because the weights change.
\end{exmp}

\begin{exmp}[Simple IPW calculation]
A treated unit has outcome 12 and propensity score 0.25. Its treated contribution
to the ATE IPW sum is
\[
    \frac{12}{0.25}=48.
\]
The large weight reflects that treated units with this \(X\) value are rare.
Large weights make estimates sensitive and should trigger overlap diagnostics.
\end{exmp}

\section{Checklist for Students}

After this lecture you should be able to:
\begin{itemize}
    \item define the conditional expectation function;
    \item use the law of iterated expectations;
    \item explain the CEF decomposition property;
    \item define population regression as a best linear predictor;
    \item interpret regression anatomy and partialling out;
    \item distinguish mechanical regression interpretation from causal interpretation;
    \item state conditional independence in potential-outcomes notation;
    \item explain how CIA identifies conditional treatment effects;
    \item distinguish ATE and ATT weighting;
    \item derive and sign the omitted-variable-bias formula;
    \item identify good controls, mediators, and bad controls;
    \item explain matching, common support, propensity scores, and IPW;
    \item explain why grouped regressions often need cell-size weights;
    \item discuss OLS, probit, and logit for limited dependent variables;
    \item explain why conditional-on-positive effects may not be causal.
\end{itemize}

\section{Practice Questions}

\begin{enumerate}
    \item Define the CEF and give an example using wages and schooling.
    \item Prove that \(Y_i-\E[Y_i\mid X_i]\) has conditional mean zero given
    \(X_i\).
    \item State the law of iterated expectations and explain it in words.
    \item Define the population regression coefficient as a minimization
    problem.
    \item Why does orthogonality of OLS residuals not imply causality?
    \item Explain regression anatomy in one sentence.
    \item What does conditional independence mean?
    \item Under CIA, why can treated and untreated units with the same \(X\) be
    compared?
    \item Give an example where ATE and ATT differ.
    \item Derive the omitted-variable-bias formula for one omitted variable.
    \item Give the sign of OVB when the omitted variable raises the outcome and
    is more common among treated units.
    \item What makes a control a bad control?
    \item Why can controlling for a mediator be wrong if the target is the total
    effect?
    \item Explain common support.
    \item State the propensity-score theorem.
    \item Why is high treatment-prediction accuracy not necessarily good for a
    propensity-score design?
    \item Write the IPW formula for the ATE and explain the weights.
    \item When should grouped regressions be weighted by cell size?
    \item Why can OLS be useful for a binary outcome?
    \item Why are conditional-on-positive outcome regressions often hard to
    interpret causally?
\end{enumerate}

\end{document}
